\DTMsavedate{mydate}{2022-06-30}
\maketitle{Modern Physics}{Applied Physics}{\DTMusedate{mydate}}

% ----------------------------------------------------- %
% COURSE INFO
\subsection*{Course Information}
\vspace{0ex}%
\noindent\parbox{0.5\textwidth}{%
    \noindent\begin{tabular}{@{}ll}
                 \textsf{Course:}          & Modern Physics \\
                 \textsf{Course Code:}         & 	SPH 1104    \\
                 \textsf{Credit Hours:}    & 	10
    \end{tabular}}
\vspace{2ex}\hrule\vspace{2ex}

% ----------------------------------------------------- %
% INSTRUCTOR INFO
\subsection*{Lecturer Information}
\noindent\parbox{0.5\textwidth}{%
    \noindent\begin{tabular}{@{}ll}
                 \textsf{Name:}         &    Innocent Muchingami           \\
% \textsf{Phone:}        &    \phone          \\
\textsf{Email:}        &    \href{mailto:innocent.muchingami@nust.ac.zw}{innocent.muchingami@nust.ac.zw}    \\
\textsf{Qualifications:}        & PhD in Groundwater Contaminant Transport (UWC), MSc in Geophysics \\
 &                                (NUST), PGD in Higher Education (NUST), BSc (Hons) in Applied Physics \\
&                                 (NUST)

    \end{tabular}}
\vspace{2ex}\hrule\vspace{2ex}


% ----------------------------------------------------- %
% COURSE SYNOPSIS
\subsection*{Course Synopsis}
The module looks at the particle nature of radiation - The photon: Planck's postulate and thermal
radiation, Blackbody radiation, the photoelectric effect, the Compton effect, X-ray production
and pair production; Interaction of radiation with matter-photon emission and absorption;
Stationery states, discrete energy spectrum and the continuous energy spectrum; The Frank-Hertz
experiment; Spontaneous and stimulated emission; The Wave nature of particles - The matter
wave: De Broglie's Postulate; The electron diffraction experiment; The wave-particle duality;
The uncertainty principle; The properties of matter waves; The Thomson and Rutherford mode;
The stability of the atom and Bohr's Postulates and his model of the atom; Atomic spectra; The
Hydrogen Atom; Correction for finite nuclear mass; The Nuclear Models: Nuclear properties,
sizes and densities, masses and densities; The Nuclear Models - Liquid drop; The deuteron; Shell
Fermi gas models; Binding energy nuclear forces; Magic numbers and the nuclear decay and
nuclear reactions, e-capture a, and emission; Fission and fusion and other nuclear reactions; The
origin of elements; Introduction to Elementary Particles: Isospin, Pions, Leptons and Families of
elementary particles.
\vspace{2ex} \hrule\vspace{2ex}
% ----------------------------------------------------- %
% COURSE TOPICS
\subsection*{Topics}
\begin{outline}
    \1
    \1
    \1
    \1
    \1
\end{outline}
\vspace{2ex} \hrule\vspace{2ex}

% ----------------------------------------------------- %
% Students Assenssment
\subsection*{Assessment of Students}
\begin{outline}
    \1 Students are assessed through written assignments and written in-class tests that shall form part of course work assessment mark.
    \1 There shall be a final examination at the end of the semester.
\end{outline}
\vspace{2ex}\hrule\vspace{2ex}

% ----------------------------------------------------- %
% Grade Evaluation
\subsection*{Course Evaluation and Grading Policies}
Grades are based on:
Continuous assessment (\qty{25}{\percent}),
Final Exam (\qty{75}{\percent})
\vspace{2ex}\hrule
\vspace{2ex}\hrule\vspace{2ex}

% ----------------------------------------------------- %
% EQUIPMENT
% https://sites.udel.edu/computing-purchases/personal-specs/
% \subsection*{Essential Equipment}

% \vspace{2ex}\hrule\vspace{2ex}

% Policies and Other information
